\documentclass[11pt]{article}
\usepackage[margin=0.88in]{geometry}
\usepackage{amsmath,amssymb,amsthm,mathtools,microtype,graphicx}
\usepackage[hidelinks]{hyperref}

\newtheorem{theorem}{Theorem}
\newtheorem{lemma}{Lemma}
\newtheorem{corollary}{Corollary}
\newcommand{\Hh}{\mathcal H}
\newcommand{\Tt}{\mathcal T}
\newcommand{\Uu}{\mathcal U}
\DeclareMathOperator{\ran}{ran}
\DeclareMathOperator{\spec}{spec}

\title{A sharp uniform noise bound for interpolating\\oblique reconstructions}
\author{Candidate full solution to the operator-norm question in arXiv:1706.06444}
\date{17 August 2026}

\begin{document}
\maketitle

\begin{center}
\fbox{\parbox{0.93\textwidth}{\textbf{Result.}
If the sampling frame has lower frame bound $A$ and
$c=\cos(\varphi_{\Tt,\Uu})>0$, then every member of the source's interpolating
family satisfies
\[
             \|Q_\lambda\|_{\mathrm{op}}
             \le \frac{1}{\sqrt A\,c},\qquad 0\le\lambda\le1.
\]
This is the best possible universal bound for every fixed choice of
$0<A\le B$, $0<c\le1$, and $\lambda$: equality is attained simultaneously
for all $\lambda$ by an explicit three-dimensional example.}}
\end{center}

\section{Source question and notation}

Berger, Gr\"ochenig, and Matz \cite{BGM} study a frame $(u_j)$ for a closed
subspace $\Uu\subset\Hh$.  Write $U:\ell^2\to\Hh$ for its synthesis
operator, $U^*$ for its analysis operator, $S=UU^*$ for its frame operator,
and $G=U^*U$ for its Gramian.  Let $A,B$ be frame bounds, so that
\[
 \spec(S|_{\Uu})\subset[A,B],\qquad 0<A\le B<\infty.
\]
For a closed reconstruction space $\Tt$, assume
\[
 c:=\cos(\varphi_{\Tt,\Uu})
   =\inf_{\substack{t\in\Tt\\\|t\|=1}}\|P_{\Uu}t\|>0.
\]
For $0<\lambda\le1$, put
\[
 \Sigma_\lambda=\lambda I+(1-\lambda)G,
 \qquad
 Q_\lambda d=\underset{g\in\Tt}{\operatorname{argmin}}
 \|\Sigma_\lambda^{-1/2}(U^*g-d)\|_2.
\]
At $\lambda=0$ the source uses $G^{\dagger/2}$ in place of
$\Sigma_\lambda^{-1/2}$.  These operators are perfect:
$Q_\lambda U^*t=t$ for $t\in\Tt$.

Lemma 3.8 of the source bounds the different quantity
$\|Q_\lambda U^*\|$.  Immediately afterward, on source PDF page 16, the
authors state that they do not know how to obtain a meaningful bound on
$\|Q_\lambda\|$ itself.  The relevant passage is reproduced below.

\begin{figure}[h]
\centering
\includegraphics[width=0.91\textwidth]{figures/open_problem_crop.png}
\caption{The extracted request after Lemma 3.8 in the source PDF, page 16.}
\end{figure}

\section{Proof intuition}

The apparent obstruction is that $\|\Sigma_\lambda^{-1/2}\|$ blows up as
$\lambda\downarrow0$ whenever the sampling frame is redundant.  But
$Q_\lambda$ annihilates the orthogonal complement of $\ran U^*$, exactly
where that blow-up occurs.  On $\ran U^*$ the preconditioner has norm
$(\lambda+(1-\lambda)A)^{-1/2}$.

The same preconditioner improves the lower bound of the weighted analysis
map on $\Tt$ by the reciprocal factor
$\sqrt{A/(\lambda+(1-\lambda)A)}$.  Taking its Moore--Penrose inverse and
multiplying the two estimates makes all $\lambda$-dependence cancel.  A
one-dimensional reconstruction space aligned with a lower-frame-bound
eigenvector makes both estimates equalities, proving optimality.

\section{The sharp theorem}

Let $J:\ell^2(I)\to\Tt$ be a unitary synthesis map for an orthonormal basis
of $\Tt$.  Define
\[
 W_\lambda=
 \begin{cases}
  (\lambda I+(1-\lambda)G)^{-1/2},&0<\lambda\le1,\\
  G^{\dagger/2},&\lambda=0,
 \end{cases}
 \qquad D_\lambda=W_\lambda U^*J.
\]

\begin{lemma}[Two compensating estimates]\label{lem:compensate}
For $d_\lambda=\lambda+(1-\lambda)A$,
\[
 \|D_\lambda x\|\ge
 c\sqrt{\frac{A}{d_\lambda}}\,\|x\|,
 \qquad
 \|W_\lambda|_{\ran U^*}\|\le\frac1{\sqrt{d_\lambda}}.
\]
Both statements include $\lambda=0$ with $d_0=A$.
\end{lemma}

\begin{proof}
Because $(u_j)$ is a frame for $\Uu$, $\ran U^*$ is closed.  Under the polar
decomposition of $U^*|_{\Uu}$, the restriction of $G$ to $\ran U^*$ is
unitarily equivalent to $S|_{\Uu}$.  Hence
\[
 \spec(G|_{\ran U^*})\subset[A,B],
\]
and functional calculus gives the second estimate.

The same polar decomposition, or the standard push-through identity for
frame operators, gives for $t\in\Tt$
\begin{align*}
 \|W_\lambda U^*t\|^2
 &=\left\langle P_{\Uu}t,
 S\bigl(\lambda I+(1-\lambda)S\bigr)^{-1}P_{\Uu}t\right\rangle.
\end{align*}
The scalar function $s\mapsto s/(\lambda+(1-\lambda)s)$ is increasing on
$(0,\infty)$ (constant when $\lambda=0$).  Spectral calculus and the angle
condition therefore yield
\[
 \|W_\lambda U^*t\|^2
 \ge \frac{A}{d_\lambda}\|P_{\Uu}t\|^2
 \ge \frac{A}{d_\lambda}c^2\|t\|^2.
\]
Taking $t=Jx$ proves the first estimate.  At $\lambda=0$ the middle operator
is $P_{\Uu}$, so the same formula remains valid.
\end{proof}

\begin{theorem}[Optimal universal noise bound]\label{thm:main}
For every $0\le\lambda\le1$,
\[
                 \boxed{\displaystyle
                 \|Q_\lambda\|_{\mathrm{op}}
                 \le\frac1{\sqrt A\,c}.}
\]
Moreover this constant is optimal even if $A,B,c$, and $\lambda$ are all
fixed in advance.
\end{theorem}

\begin{proof}
The representation theorem for $Q_\lambda$ in the source, applied to an
orthonormal basis of $\Tt$, reads
\[
                    Q_\lambda=J D_\lambda^\dagger W_\lambda.
\]
Here $D_\lambda$ has closed range by Lemma \ref{lem:compensate}.  Its range
is contained in $\ran U^*$, so $D_\lambda^\dagger$ annihilates
$(\ran U^*)^\perp$.  Since $W_\lambda$ preserves the orthogonal decomposition
$\ell^2=\ran U^*\oplus(\ran U^*)^\perp$, Lemma \ref{lem:compensate} gives
\begin{align*}
 \|Q_\lambda\|
 &\le \|D_\lambda^\dagger\|
       \|W_\lambda|_{\ran U^*}\|\\
 &\le \frac1c\sqrt{\frac{d_\lambda}{A}}
       \frac1{\sqrt{d_\lambda}}
  =\frac1{\sqrt A\,c}.
\end{align*}

For sharpness, fix arbitrary $0<A\le B$ and $0<c\le1$.  In
$\Hh=\mathbb C^3$ with orthonormal basis $e_1,e_2,e_3$, take
\[
 \Uu=\operatorname{span}\{e_1,e_2\},\quad
 u_1=\sqrt A\,e_1,\quad u_2=\sqrt B\,e_2,
 \quad
 \Tt=\operatorname{span}\{t\},
 \quad t=ce_1+\sqrt{1-c^2}\,e_3.
\]
The sampling frame has optimal bounds $A,B$, and the subspace angle has
cosine $c$.  If $d=(d_1,d_2)$ and $g=\alpha t$, the weighted least-squares
functional is
\[
 \frac{|\sqrt A\,c\alpha-d_1|^2}{\lambda+(1-\lambda)A}
 +\frac{|d_2|^2}{\lambda+(1-\lambda)B}.
\]
Its unique minimizer is
$\alpha=d_1/(\sqrt A\,c)$, including the Moore--Penrose endpoint
$\lambda=0$.  Thus
\[
 Q_\lambda(d_1,d_2)=\frac{d_1}{\sqrt A\,c}\,t,
 \qquad \|Q_\lambda\|=\frac1{\sqrt A\,c}
\]
for every $\lambda\in[0,1]$.  This attains the asserted bound with all
parameters fixed.
\end{proof}

\begin{corollary}[Complete universal bracket]
Put
\[
 a_{\Tt}:=\inf_{\substack{t\in\Tt\\\|t\|=1}}\|U^*t\|>0.
\]
Then
\[
 \frac1{a_{\Tt}}=\|Q_1\|\le\|Q_\lambda\|
 \le\frac1{\sqrt A\,c},\qquad 0\le\lambda\le1.
\]
Both inequalities can be equalities simultaneously.
\end{corollary}

\begin{proof}
$Q_1=(U^*|_{\Tt})^\dagger$ is the Moore--Penrose left inverse, so its norm
is $1/a_{\Tt}$.  Every $Q_\lambda$ is a left inverse of $U^*|_{\Tt}$;
testing on unit vectors approaching the lower modulus $a_{\Tt}$ shows that
no left inverse has smaller norm.  The upper estimate is Theorem
\ref{thm:main}, and the sharpness example has
$a_{\Tt}=\sqrt A\,c$.
\end{proof}

\section{Verification, novelty, and scope}

The included script constructs random finite frames, redundant analysis
operators, reconstruction subspaces, and all endpoint/interior values of
$\lambda$.  It checks perfect reconstruction and the theorem's inequality.
It separately verifies the equality construction for several prescribed
values of $A,B,c$.  These computations are sanity checks only; the proof is
the functional-calculus argument above.

A bounded novelty search on 17 August 2026 covered the run's four lightweight
indexes and web searches for the exact title together with the phrase
\emph{operator norm $Q_\lambda$}, the source's phrase
\emph{meaningful bound}, and weighted pseudoinverse norm bounds.  The exact
hits were the arXiv and published versions of the source; no later paper
explicitly resolving this request was found.  Novelty is therefore plausible,
not certified.

The result completely determines the best universal bound in terms of the
source's frame bounds and subspace angle.  It does not claim an exact formula
for a specified frame, nor does it prove or disprove monotonicity of
$\lambda\mapsto\|Q_\lambda\|$ for every individual frame.  Human review
should focus on the range restriction before the norm estimate and the
$\lambda=0$ Moore--Penrose endpoint.  Verdict:
\textbf{candidate full solution, likely valid}.

\begin{thebibliography}{9}
\bibitem{BGM}
P. Berger, K. Gr\"ochenig, and G. Matz,
\emph{Sampling and Reconstruction in Distinct Subspaces Using Oblique
Projections}, arXiv:1706.06444; J. Fourier Anal. Appl. 25 (2019), 1080--1112.
\end{thebibliography}

\end{document}
